% Responsável: TODO(grupo) — frente "Dados e protocolo"
% Enunciado, item 6.3: origem, setup físico, taxa de amostragem, tipos de falha
% rotulados, quantidade de amostras por classe.
\chapter{Descrição do dataset}
\label{cap:dataset}

\section{Origem e setup experimental}
\label{sec:setup}

TODO(grupo): bancada do \emph{Case Western Reserve University Bearing Data
Center}; motor, dinamômetro, posição dos acelerômetros, defeitos induzidos por
eletroerosão.

\section{Configuração adotada}
\label{sec:config-adotada}

% Todos os números desta tabela devem ser conferidos contra a documentação
% oficial do CWRU antes da entrega. Não copiar de fonte secundária.
\begin{table}[htb]
  \centering
  \caption{Configuração do subconjunto do CWRU utilizada neste trabalho.}
  \label{tab:config}
  \begin{tabular}{ll}
    \toprule
    \textbf{Parâmetro} & \textbf{Valor} \\
    \midrule
    Sensor                & Drive end (DE) \\
    Taxa de amostragem    & \SI{12}{\kilo\hertz} \\
    Rolamento             & SKF 6205-2RS JEM \\
    Classes               & normal, pista interna, pista externa, esfera \\
    Diâmetros de defeito  & \SI{0.007}{\inch}, \SI{0.014}{\inch}, \SI{0.021}{\inch} \\
    Cargas                & \SIrange{0}{3}{\hp} \\
    Posição do defeito em pista externa & TODO(grupo): declarar 3, 6 ou 12 horas \\
    \bottomrule
  \end{tabular}
  \fonte{Elaborado pelos autores a partir da documentação do CWRU Bearing Data Center.}
\end{table}

Justificativa das escolhas: TODO(grupo).

\section{Quantidade de amostras por classe}
\label{sec:amostras-por-classe}

% Preencher SOMENTE com a contagem produzida por src/data.py. Nada de estimativa.
TODO(grupo): tabela de janelas por classe e por carga, após a segmentação.

\section{Auditoria dos arquivos}
\label{sec:auditoria}

Nem toda gravação do CWRU é utilizável. TODO(grupo): resumir os critérios de
exclusão adotados; a lista completa está no Apêndice~\ref{ap:arquivos-excluidos}.
